% ===============================================================
% GRADE 9 -- MATATAG Science 9, Three-Term BOW (rev. 17 Apr 2026)
% ===============================================================
\section{Grade 9}
\resetcards
\begin{gradebody}

% ---------------------------------------------------------------
\subsection{Term 1: Newton's Laws, Electricity, EM Waves, Plates}

\topic{Newton's Laws of Motion}
\kf
\begin{keyfacts}
\item \textbf{1st law (inertia):} an object stays at rest or in uniform motion unless an unbalanced (net) force acts on it. More mass means more inertia.
\item \textbf{2nd law:} $F = ma$. Acceleration increases with net force and decreases with mass.
\item \textbf{3rd law:} for every action there is an equal and opposite reaction. The two forces act on \emph{different} objects, so they do not cancel.
\item $\SI{1}{N} = \SI{1}{\kilogram\meter\per\second\squared}$. The laws were published in the \emph{Principia} (1687).
\end{keyfacts}
\cards
\qa{Newton's law of inertia}{First law}{E}
\qa{Tendency to resist a change in motion}{Inertia}{E}
\qa{Quantity that measures inertia}{Mass}{A}
\qa{Passengers lurch forward when a jeepney brakes}{First law (inertia)}{E}
\qa{Coin drops into the glass when the card is flicked}{Inertia}{A}
\qa{$F = ma$ is which law?}{Second law}{E}
\qa{Net force on a \SI{2}{kg} cart accelerating at \SI{3}{\meter\per\second\squared}}{\SI{6}{N}}{A}
\qa{Same force, double the mass: acceleration}{Halved}{D}
\qa{Swimmer pushes water back and moves forward}{Third law}{E}
\qa{Stepping off a bangka pushes the boat}{Backward (away from you)}{A}
\qa{Why action and reaction don't cancel}{They act on different objects}{D}
\qa{Book on a table: reaction to Earth pulling the book down}{Book pulls Earth up}{D}
\qa{1 newton in base units}{\si{\kilogram\meter\per\second\squared}}{D}
\qa{Newton's 1687 book}{\emph{Principia}}{D}
\begin{confbox}
\confusion{Action--reaction pair}{balanced forces}{an action--reaction pair acts on two objects; balanced forces act on one object.}
\confusion{Mass}{inertia}{mass is the \emph{measure} of inertia; inertia is the \emph{property}.}
\end{confbox}

\topic{Electricity and Circuits}
\kf
\begin{keyfacts}
\item Current is a flow of electrons, measured in amperes (A) with an \textbf{ammeter in series}.
\item Voltage (potential difference) is measured in volts (V) with a \textbf{voltmeter in parallel}. Resistance is in ohms (\si{\ohm}).
\item \textbf{Ohm's law} (Georg Ohm): $V = IR$.
\item \textbf{Series:} one path, same current everywhere, voltages add, $R_T = R_1 + R_2$. One bulb out turns all off, and more bulbs make them dimmer.
\item \textbf{Parallel:} many paths, same voltage across each branch, currents add. One bulb out leaves the others on. Homes are wired in parallel.
\item Electrons flow from $-$ to $+$; conventional current flows from $+$ to $-$.
\item Safety: fuses and circuit breakers, no overloaded outlets, dry hands. PH household supply is \SI{220}{V}.
\end{keyfacts}
\cards
\qa{Unit of current}{Ampere}{E}
\qa{Unit of resistance}{Ohm}{E}
\qa{Unit of potential difference}{Volt}{E}
\qa{Meter connected in series}{Ammeter}{A}
\qa{Meter connected in parallel}{Voltmeter}{A}
\qa{$V = IR$}{Ohm's law}{E}
\qa{\SI{12}{V} across \SI{4}{\ohm}: current}{\SI{3}{A}}{A}
\qa{\SI{2}{\ohm} and \SI{3}{\ohm} in series}{\SI{5}{\ohm}}{A}
\qa{One bulb blows and all go off}{Series circuit}{E}
\qa{Type of household wiring}{Parallel}{E}
\qa{Adding bulbs in series makes them}{Dimmer}{A}
\qa{Same across every parallel branch}{Voltage}{D}
\qa{Same everywhere in a series circuit}{Current}{D}
\qa{Device that cuts the circuit on overload}{Fuse / circuit breaker}{A}
\qa{Standard household voltage in PH}{\SI{220}{V}}{A}
\qa{Direction of electron flow}{Negative to positive terminal}{D}
\begin{confbox}
\confusion{Current}{voltage}{current is the flow (A); voltage is the push (V).}
\confusion{Series}{parallel}{series has one path and shares voltage; parallel has many paths and each branch gets full voltage.}
\confusion{Fuse}{circuit breaker}{a fuse melts and must be replaced; a breaker trips and can be reset.}
\end{confbox}
\mnemonic{Ohm's triangle: \textbf{V} on top, \textbf{I} and \textbf{R} below. Cover the one you need.}

\topic{Electromagnetic Waves}
\kf
\begin{keyfacts}
\item EMR is energy from vibrating charged particles. It travels as \textbf{transverse} waves, needs no medium, and moves at the speed of light in a vacuum.
\item Order of \emph{increasing} frequency and energy (decreasing wavelength): radio, microwave, infrared, visible, ultraviolet, X-ray, gamma.
\item $v = f\lambda$. Frequency is in hertz (Hz).
\item Uses: radio/TV and telecoms (radio), ovens, Wi-Fi and satellites (microwave), remotes and thermal imaging (IR), sterilising and spotting fake bills (UV), bone imaging and airport scanners (X-ray), cancer therapy and sterilising equipment (gamma).
\item Hazards: UV causes sunburn, skin cancer, and cataracts. X-rays and gamma are ionizing and damage cells and DNA.
\item Discoverers: \textbf{Maxwell} predicted EM waves, \textbf{Hertz} produced radio waves, \textbf{Herschel} found IR, \textbf{R\"ontgen} found X-rays (1895).
\end{keyfacts}
\cards
\qa{EM waves are transverse or longitudinal?}{Transverse}{E}
\qa{Lowest-frequency EM wave}{Radio}{E}
\qa{Highest-energy EM wave}{Gamma}{E}
\qa{EM wave used by TV remotes}{Infrared}{E}
\qa{Causes sunburn and is used for sterilizing}{Ultraviolet}{E}
\qa{EM wave for imaging bones}{X-ray}{E}
\qa{Used in cancer radiotherapy}{Gamma rays}{A}
\qa{Used by microwave ovens and Wi-Fi}{Microwaves}{A}
\qa{Distance from crest to crest}{Wavelength}{E}
\qa{Unit of frequency}{Hertz}{E}
\qa{As frequency rises, wavelength}{Decreases}{A}
\qa{Visible colour with the longest wavelength}{Red}{A}
\qa{$f = \SI{100}{Hz}$, $\lambda = \SI{2}{m}$: wave speed}{\SI{200}{\meter\per\second}}{A}
\qa{Discovered X-rays}{Wilhelm R\"ontgen}{A}
\qa{Predicted EM waves mathematically}{James Clerk Maxwell}{D}
\qa{First produced and detected radio waves}{Heinrich Hertz}{D}
\qa{Discovered infrared}{William Herschel}{D}
\qa{Do X-rays travel faster than radio waves in a vacuum?}{No, both travel at $c$}{D}
\qa{Ionizing parts of the EM spectrum}{High-energy UV, X-ray, gamma}{D}
\begin{confbox}
\confusion{Frequency}{wavelength}{they are inversely related: high frequency means short wavelength.}
\confusion{Amplitude}{wavelength}{amplitude is the wave's height (intensity); wavelength is the crest-to-crest length.}
\confusion{Ionizing}{non-ionizing}{ionizing radiation (X-ray, gamma) can knock electrons off atoms and damage DNA; radio, microwave, and IR cannot.}
\end{confbox}
\mnemonic{\textbf{R}aging \textbf{M}artians \textbf{I}nvaded \textbf{V}enus \textbf{U}sing \textbf{X}-ray \textbf{G}uns (radio $\rightarrow$ gamma).}

\topic{Continental Drift and Plate Boundaries}
\kf
\begin{keyfacts}
\item \textbf{Alfred Wegener} proposed continental drift: one supercontinent, \textbf{Pangaea}, surrounded by Panthalassa.
\item Evidence: jigsaw fit of coastlines (South America and Africa), matching rocks and mountains, fossils (\emph{Mesosaurus}, \emph{Glossopteris}), and old glacier deposits.
\item Wegener's idea was rejected at first because he had no mechanism. \textbf{Harry Hess} later proposed seafloor spreading.
\item \textbf{Divergent} boundaries pull apart (mid-ocean ridges, rift valleys: Mid-Atlantic Ridge, East African Rift).
\item \textbf{Convergent} boundaries collide: ocean--ocean gives an island arc and trench; ocean--continent gives a trench and volcanic mountains (Andes); continent--continent gives folded mountains (Himalayas).
\item \textbf{Transform} boundaries slide past each other (San Andreas Fault).
\item The PH sits between the \textbf{Philippine Sea Plate} (subducting at the Philippine Trench, east) and the Eurasian/Sunda Plate (subducting at the \textbf{Manila Trench}, west).
\end{keyfacts}
\cards
\qa{Proposed continental drift}{Alfred Wegener}{E}
\qa{Ancient supercontinent}{Pangaea}{E}
\qa{Ocean surrounding Pangaea}{Panthalassa}{D}
\qa{Freshwater reptile fossil in both South America and Africa}{\emph{Mesosaurus}}{A}
\qa{Seed-fern fossil used as drift evidence}{\emph{Glossopteris}}{A}
\qa{Why drift was rejected at first}{No mechanism explained}{D}
\qa{Proposed seafloor spreading}{Harry Hess}{D}
\qa{Boundary where plates move apart}{Divergent}{E}
\qa{Boundary where plates slide past each other}{Transform}{E}
\qa{Example of a transform boundary}{San Andreas Fault}{A}
\qa{Example of a divergent boundary}{Mid-Atlantic Ridge}{A}
\qa{Ocean--continent convergence forms}{Trench + volcanic mountains}{A}
\qa{Continent--continent convergence forms}{Folded mountains (Himalayas)}{A}
\qa{Ocean--ocean convergence forms}{Volcanic island arc}{D}
\qa{Plate that subducts at the Philippine Trench}{Philippine Sea Plate}{D}
\qa{Trench west of Luzon}{Manila Trench}{D}
\begin{confbox}
\confusion{Continental drift}{plate tectonics}{drift says continents move; plate tectonics explains \emph{how} (lithospheric plates on the asthenosphere).}
\confusion{Trench}{ridge}{a trench forms where plates converge (subduction); a ridge forms where they diverge.}
\end{confbox}

% ---------------------------------------------------------------
\subsection{Term 2: Geologic Time, Earth's Interior, Solar System, Biodiversity}

\topic{Fossils, Dating, and the Geologic Time Scale}
\kf
\begin{keyfacts}
\item \textbf{Relative dating} gives the order of events. By the \textbf{law of superposition} (Nicolas Steno), the oldest undisturbed layer is at the bottom. \textbf{Index fossils} are widespread but lived for a short time.
\item \textbf{Absolute (radiometric) dating} gives an age in years from radioactive decay and half-life. Carbon-14 (half-life about 5\,730 years) dates once-living material; uranium and potassium methods date rocks.
\item Time units: eon $>$ era $>$ period $>$ epoch. Earth is about 4.6 billion years old.
\item Eras: Paleozoic (``ancient life''), Mesozoic (age of reptiles/dinosaurs), Cenozoic (age of mammals; today). We live in the Holocene epoch.
\item The largest mass extinction was at the end of the Permian. The dinosaurs died at the end of the Cretaceous.
\end{keyfacts}
\cards
\qa{Oldest undisturbed layer is at the bottom}{Law of superposition}{E}
\qa{Formulated the law of superposition}{Nicolas Steno}{D}
\qa{Fossils best for matching rock ages}{Index fossils}{A}
\qa{Dating that gives an age in years}{Absolute (radiometric)}{E}
\qa{Time for half of a radioactive sample to decay}{Half-life}{E}
\qa{Half-life of carbon-14}{About 5\,730 years}{D}
\qa{Approximate age of Earth}{4.6 billion years}{E}
\qa{Largest division of geologic time}{Eon}{A}
\qa{``Age of reptiles''}{Mesozoic Era}{E}
\qa{Era we live in}{Cenozoic}{A}
\qa{Current epoch}{Holocene}{D}
\qa{Largest mass extinction}{End-Permian}{D}
\begin{confbox}
\confusion{Relative}{absolute dating}{relative gives older or younger; absolute gives a number of years.}
\confusion{Eon}{era}{an eon is larger; eras are subdivisions of an eon.}
\end{confbox}

\topic{Earth's Interior and Seismic Waves}
\kf
\begin{keyfacts}
\item Layers: crust (thinnest), mantle (thickest), liquid outer core, solid inner core (iron--nickel). The lithosphere is rigid; the asthenosphere below it is plastic, and the plates move on it.
\item \textbf{P-waves} are primary, compressional, and fastest, and pass through solids and liquids. \textbf{S-waves} are secondary and transverse and pass through solids only. \textbf{Surface waves} are the most damaging.
\item The S-wave shadow zone shows that the outer core is liquid.
\item Boundaries: \textbf{Moho} (crust--mantle, Mohorovi\v{c}i\'c), \textbf{Gutenberg} (mantle--core). \textbf{Inge Lehmann} discovered the inner core.
\item Earth's internal heat comes mainly from radioactive decay. The moving outer core generates the magnetic field.
\end{keyfacts}
\cards
\qa{Fastest seismic wave}{P-wave}{E}
\qa{Seismic wave that cannot pass through liquid}{S-wave}{A}
\qa{Evidence that the outer core is liquid}{S-wave shadow zone}{D}
\qa{Crust--mantle boundary}{Moho}{A}
\qa{Discovered the inner core}{Inge Lehmann}{D}
\qa{Thickest layer}{Mantle}{E}
\qa{Thinnest layer}{Crust}{E}
\qa{Plastic layer the plates ride on}{Asthenosphere}{A}
\qa{Main metals of the core}{Iron and nickel}{E}
\qa{Layer that generates Earth's magnetic field}{Outer core}{A}
\qa{Main source of Earth's internal heat}{Radioactive decay}{A}
\begin{confbox}
\confusion{P-wave}{S-wave}{P waves are push--pull and travel through anything; S waves shake side to side through solids only.}
\confusion{Lithosphere}{asthenosphere}{the lithosphere is rigid (crust + top mantle); the asthenosphere is soft and flowing.}
\confusion{Crust}{lithosphere}{the lithosphere includes the crust \emph{plus} the uppermost mantle.}
\end{confbox}
\mnemonic{\textbf{P} = \textbf{P}rimary, \textbf{P}ush--pull, \textbf{P}asses all. \textbf{S} = \textbf{S}econdary, \textbf{S}hake, \textbf{S}olids only.}

\topic{The Solar System, Small Bodies, and Space Technology}
\kf
\begin{keyfacts}
\item The \textbf{nebular hypothesis} says the Solar System formed about 4.6 billion years ago from a spinning cloud of gas and dust.
\item In 2006 the IAU made Pluto a dwarf planet. The dwarf planets are Ceres, Pluto, Eris, Haumea, and Makemake.
\item \textbf{Asteroids} are rocky and mostly lie between Mars and Jupiter. \textbf{Comets} are icy, with tails pointing away from the Sun. Short-period comets come from the Kuiper belt, long-period ones from the Oort cloud.
\item A meteoroid is in space, a meteor is the streak of light in the sky, and a meteorite is what lands.
\item Meteor showers happen when Earth crosses comet debris: Perseids (Aug), Orionids (Oct), Leonids (Nov), Geminids (Dec, from asteroid Phaethon).
\item Technology: Galileo first used a telescope for astronomy. Telescopes are refracting (lenses) or reflecting (mirrors). Hubble launched in 1990; JWST (infrared) in 2021. Probes: Voyager 1, New Horizons (Pluto, 2015).
\item PH space efforts: \textbf{PhilSA}; \textbf{Diwata-1} was the first PH microsatellite (2016).
\end{keyfacts}
\cards
\qa{Theory that the Solar System formed from a spinning gas cloud}{Nebular hypothesis}{E}
\qa{Year Pluto was reclassified}{2006}{A}
\qa{Body that reclassified Pluto}{IAU}{A}
\qa{Dwarf planet in the asteroid belt}{Ceres}{A}
\qa{Location of the asteroid belt}{Between Mars and Jupiter}{E}
\qa{Icy body with a tail away from the Sun}{Comet}{E}
\qa{Streak of light in the sky}{Meteor}{E}
\qa{Space rock that reaches the ground}{Meteorite}{E}
\qa{Cause of meteor showers}{Earth passing through comet debris}{A}
\qa{August meteor shower}{Perseids}{A}
\qa{December shower from asteroid Phaethon}{Geminids}{D}
\qa{Orbital period of Halley's Comet}{About 76 years}{A}
\qa{Source region of long-period comets}{Oort cloud}{D}
\qa{Region beyond Neptune that includes Pluto}{Kuiper belt}{A}
\qa{Largest planet}{Jupiter}{E}
\qa{Hottest planet}{Venus}{A}
\qa{Telescope that uses mirrors}{Reflecting telescope}{A}
\qa{First to observe Jupiter's moons with a telescope}{Galileo Galilei}{E}
\qa{Infrared space telescope launched in 2021}{James Webb (JWST)}{A}
\qa{First spacecraft to enter interstellar space}{Voyager 1}{D}
\qa{Probe that flew past Pluto in 2015}{New Horizons}{D}
\qa{Philippine Space Agency}{PhilSA}{A}
\qa{First Philippine microsatellite}{Diwata-1}{D}
\begin{confbox}
\confusion{Meteoroid}{meteor}{meteoroid (in space) $\rightarrow$ meteor (light in the sky) $\rightarrow$ meteorite (on the ground).}
\confusion{Asteroid}{comet}{asteroids are rocky with no tail; comets are icy and grow tails near the Sun.}
\confusion{Refracting}{reflecting telescope}{refracting telescopes use lenses; reflecting telescopes use mirrors (all large modern ones).}
\end{confbox}
\mnemonic{\textbf{M}y \textbf{V}ery \textbf{E}ducated \textbf{M}other \textbf{J}ust \textbf{S}erved \textbf{U}s \textbf{N}achos (Mercury $\rightarrow$ Neptune).}

\topic{Biodiversity and Endangered Species}
\kf
\begin{keyfacts}
\item \textbf{Biodiversity} has three levels: genetic, species, and ecosystem. High biodiversity gives a more stable, resilient ecosystem with more alternative food sources and pathways.
\item IUCN Red List: EX, EW, \textbf{CR}, \textbf{EN}, \textbf{VU}, NT, LC (plus DD). ``Threatened'' means CR, EN, or VU.
\item Critically endangered PH species include the Philippine eagle, the tamaraw (Mindoro), the Philippine cockatoo (\emph{katala}), and the Philippine pangolin \verify{(CR)}.
\item Main threats (HIPPO): habitat loss, invasive species (janitor fish, golden apple snail), pollution, human population, overharvesting.
\item Protection: in-situ (protected areas such as Tubbataha Reefs Natural Park) and ex-situ (zoos, seed banks). Laws: NIPAS Act (RA 7586) and the Wildlife Act (RA 9147).
\end{keyfacts}
\cards
\qa{Variety of life in an area}{Biodiversity}{E}
\qa{Three levels of biodiversity}{Genetic, species, ecosystem}{A}
\qa{Species found naturally in only one place}{Endemic species}{E}
\qa{IUCN category one step riskier than Endangered}{Critically Endangered}{A}
\qa{IUCN categories that count as ``threatened''}{CR, EN, VU}{D}
\qa{IUCN status of the Philippine eagle}{Critically Endangered}{A}
\qa{Dwarf buffalo endemic to Mindoro}{Tamaraw}{E}
\qa{Local name of the Philippine cockatoo}{Katala}{A}
\qa{Local name of the whale shark}{Butanding}{E}
\qa{Leading cause of biodiversity loss}{Habitat loss}{A}
\qa{Invasive fish in Laguna de Bay and Marikina River}{Janitor fish}{A}
\qa{Invasive snail that pests rice fields}{Golden apple snail (golden kuhol)}{A}
\qa{Conservation within the natural habitat}{In-situ}{D}
\qa{Conservation outside the habitat (zoos, seed banks)}{Ex-situ}{D}
\qa{UNESCO reef park in the Sulu Sea}{Tubbataha Reefs}{A}
\qa{Law creating the protected-areas system}{NIPAS Act (RA 7586)}{D}
\qa{PH Wildlife Act}{RA 9147}{D}
\qa{Why high biodiversity keeps ecosystems stable}{Alternative food sources and pathways}{D}
\begin{confbox}
\confusion{Extinct}{extinct in the wild}{extinct means no individuals left anywhere; EW means they survive only in captivity.}
\confusion{Endemic}{native}{endemic means found \emph{only} there; native means naturally there but possibly elsewhere too.}
\confusion{In-situ}{ex-situ}{in-situ is in the natural habitat; ex-situ is outside it.}
\end{confbox}
\mnemonic{\textbf{HIPPO}: Habitat loss, Invasive species, Pollution, Population (human), Overharvesting.}

\topic{Philippine Ecosystems and Human Impact}
\kf
\begin{keyfacts}
\item \textbf{Tropical rainforest}: high rainfall and biodiversity, with layers (emergent, canopy, understory, floor).
\item \textbf{Mangroves} are salt-tolerant intertidal trees with prop roots and pneumatophores (breathing roots). They are fish nurseries and buffer coasts against surges.
\item \textbf{Estuary}: where a river meets the sea, with brackish and very productive water.
\item \textbf{Swamp}: a tree-dominated wetland (e.g., Agusan Marsh).
\item \textbf{Coral reefs} are built by coral polyps (animals) that host algae (zooxanthellae). Heat stress causes bleaching. The PH is in the \textbf{Coral Triangle}.
\item Human threats: deforestation and \emph{kaingin}, pollution, dynamite and cyanide fishing, invasive species, and converting mangroves to fishponds.
\end{keyfacts}
\cards
\qa{Mix of fresh and salt water}{Brackish}{E}
\qa{Where a river meets the sea}{Estuary}{E}
\qa{Breathing roots of mangroves}{Pneumatophores}{D}
\qa{Local term for slash-and-burn farming}{Kaingin}{E}
\qa{Animals that build coral reefs}{Coral polyps}{A}
\qa{Algae living inside coral tissue}{Zooxanthellae}{D}
\qa{Corals turning white from heat stress}{Coral bleaching}{A}
\qa{Marine biodiversity hotspot that includes the PH}{Coral Triangle}{A}
\qa{Large wetland in Agusan del Sur}{Agusan Marsh}{A}
\qa{Coastal protection role of mangroves}{Buffer against surges and erosion}{A}
\qa{Dense ``roof'' layer of a rainforest}{Canopy}{E}
\qa{Two destructive fishing methods}{Dynamite and cyanide fishing}{E}
\qa{Example of an abiotic factor}{Sunlight, water, temperature}{E}
\qa{Major historic cause of PH mangrove loss}{\verify{Conversion to fishponds}}{D}
\begin{confbox}
\confusion{Biotic}{abiotic}{biotic factors are living (plants, animals); abiotic are nonliving (light, water, soil).}
\confusion{Swamp}{marsh}{swamps are dominated by trees; marshes by grasses and reeds.}
\confusion{Estuary}{mangrove}{an estuary is a \emph{place} (river mouth); a mangrove is a \emph{forest type} often found there.}
\end{confbox}

% ---------------------------------------------------------------
\subsection{Term 3: DNA and Mutations, Chemical Bonding}

\topic{DNA, Genes, Chromosomes, and Mutations}
\kf
\begin{keyfacts}
\item DNA (deoxyribonucleic acid) is a \textbf{double helix} (Watson and Crick, 1953), built using \textbf{Rosalind Franklin}'s X-ray image (Photo 51).
\item A nucleotide is phosphate + deoxyribose + base. Bases pair \textbf{A--T} and \textbf{G--C}. The sugar--phosphate backbone forms the sides; the base pairs form the rungs.
\item A gene is a DNA segment that codes for a protein or trait. A chromosome is coiled DNA. Humans have 23 pairs (22 autosome pairs + XX or XY).
\item Replication is semi-conservative. Helicase unzips the DNA and DNA polymerase builds the new strand.
\item \textbf{Mutation} is a change in DNA or chromosomes. Gene mutations: substitution, insertion, deletion (frameshift). Chromosomal mutations: deletion, duplication, inversion, translocation, and extra or missing chromosomes (trisomy 21 is Down syndrome).
\item Mutagens: radiation (UV, X-ray), chemicals (tobacco smoke), and infectious agents (some viruses). Effects can be harmful, beneficial, or neutral.
\end{keyfacts}
\cards
\qa{DNA stands for}{Deoxyribonucleic acid}{E}
\qa{Shape of DNA}{Double helix}{E}
\qa{Proposed the double-helix model (1953)}{Watson and Crick}{E}
\qa{Scientist behind Photo 51}{Rosalind Franklin}{A}
\qa{Base that pairs with adenine}{Thymine}{E}
\qa{Base that pairs with guanine}{Cytosine}{E}
\qa{Sugar in DNA}{Deoxyribose}{A}
\qa{Three parts of a nucleotide}{Phosphate, sugar, base}{A}
\qa{DNA segment coding for a trait}{Gene}{E}
\qa{Chromosome pairs in a human body cell}{23}{E}
\qa{Sex chromosomes of a male}{XY}{E}
\qa{Enzyme that unzips DNA}{Helicase}{D}
\qa{Each new DNA keeps one old strand}{Semi-conservative replication}{D}
\qa{Change in the DNA sequence}{Mutation}{E}
\qa{Agent that causes mutation}{Mutagen}{A}
\qa{Extra copy of chromosome 21}{Down syndrome}{A}
\qa{Mutation with no effect}{Neutral (silent)}{A}
\qa{Insertion or deletion that shifts the reading frame}{Frameshift mutation}{D}
\qa{Disorder from one base substitution in the haemoglobin gene}{Sickle-cell anaemia}{D}
\qa{If 20\% of the bases are A, the percentage of G}{30\%}{D}
\begin{confbox}
\confusion{DNA}{gene}{a gene is one short section of a DNA molecule; a chromosome is one long DNA molecule packed with proteins.}
\confusion{Gene mutation}{chromosomal mutation}{a gene mutation changes a few bases; a chromosomal mutation changes large segments or whole chromosomes.}
\confusion{DNA}{RNA}{DNA has deoxyribose, T, and two strands; RNA has ribose, U, and one strand.}
\end{confbox}

\topic{Chemical Change and Valence Electrons}
\kf
\begin{keyfacts}
\item A chemical change breaks old bonds and forms new ones, making a new substance.
\item Evidence tests: \ce{CO2} turns \textbf{limewater} milky (a \ce{CaCO3} precipitate). \ce{AgNO3} + salt solution gives a white \ce{AgCl} precipitate.
\item \textbf{Valence electrons} are outer-shell electrons and take part in bonding. Oxygen (Group 16) has 6.
\item The \textbf{octet rule}: atoms tend to gain, lose, or share electrons to have 8 valence electrons. Lewis dot diagrams show valence electrons.
\end{keyfacts}
\cards
\qa{Electrons in the outermost shell}{Valence electrons}{E}
\qa{Valence electrons of oxygen}{6}{E}
\qa{Atoms tend to reach 8 valence electrons}{Octet rule}{E}
\qa{Electron-dot diagrams are named after}{Gilbert N.\ Lewis}{A}
\qa{Test for \ce{CO2}}{Limewater turns milky}{A}
\qa{White precipitate from silver nitrate + salt water}{Silver chloride}{D}

\topic{Ionic, Covalent, and Metallic Bonding}
\kf
\begin{keyfacts}
\item \textbf{Ionic} bonds form when a metal transfers electrons to a nonmetal: \ce{Na -> Na+ + e-} and \ce{Cl + e- -> Cl-}. The ions form a \textbf{crystal lattice} that has a high melting point, is brittle, and conducts when molten or dissolved.
\item Formulas: \ce{NaCl} sodium chloride, \ce{MgO} magnesium oxide, \ce{KCl} potassium chloride, \ce{MgCl2} magnesium chloride.
\item \textbf{Covalent} bonds form when nonmetals share electrons, making molecules with low melting points that are poor conductors: \ce{H2O}, \ce{CO2} (O=C=O), \ce{NH3}.
\item \textbf{Metallic} bonds hold metal cations in a \textbf{sea of electrons}, so metals conduct and are malleable, ductile, and shiny.
\end{keyfacts}
\cards
\qa{Positively charged ion}{Cation}{E}
\qa{Negatively charged ion}{Anion}{E}
\qa{How sodium becomes \ce{Na+}}{Loses 1 electron}{A}
\qa{Ion charge of magnesium}{2+}{A}
\qa{Bond formed by electron transfer}{Ionic}{E}
\qa{Bond formed by sharing electrons}{Covalent}{E}
\qa{Bond within metals}{Metallic}{E}
\qa{Model of metallic bonding}{Sea of electrons}{A}
\qa{Formula of magnesium chloride}{\ce{MgCl2}}{A}
\qa{Name of \ce{MgO}}{Magnesium oxide}{E}
\qa{Formula of ammonia}{\ce{NH3}}{A}
\qa{Structure of ionic compounds}{Crystal lattice}{A}
\qa{High melting point, brittle, conducts when dissolved}{Ionic compound}{A}
\qa{Bonds in a \ce{CO2} molecule}{Two double covalent bonds}{D}
\qa{Why molten NaCl conducts electricity}{Its ions are free to move}{D}
\qa{Why metals can be hammered into sheets}{Cations slide within the electron sea}{D}
\qa{Bonding in diamond}{Covalent (network)}{D}
\begin{confbox}
\confusion{Ionic}{covalent}{ionic transfers electrons (metal + nonmetal); covalent shares them (nonmetal + nonmetal).}
\confusion{Cation}{anion}{a cation \emph{lost} electrons ($+$); an anion \emph{gained} electrons ($-$).}
\confusion{Molecule}{crystal lattice}{covalent compounds form separate molecules; ionic compounds form a continuous lattice (no single ``molecule'').}
\end{confbox}
\mnemonic{Ca\textbf{t}ions are ``pawsi\textbf{t}ive'' (the $t$ looks like $+$); \textbf{A}nion = \textbf{A} \textbf{N}egative \textbf{ION}.}

\end{gradebody}
\cardstats
